\subsubsection{Closed-Form Computation Using Only \texorpdfstring{$\Delta$}{Delta}}\label{app:delta-only}
Let $z_\star(u)=\lambda(u)^{-1}$ be the smallest positive root of $\Delta(z,u)$, and define
\[
C(u):=\lim_{z\to z_\star}(1-z/z_\star)\zeta(z,u)=\frac{1}{-z_\star\,\partial_z\Delta(z_\star,u)}.
\]
By Witt inversion and $\zeta=1/\Delta$, one obtains
\[
\log\mathfrak{M}(\theta)
=\log C(u)+\sum_{m\ge 2}\frac{\mu(m)}{m}\log\zeta(z_\star^m,u^m)
=\log C(u)-\sum_{m\ge 2}\frac{\mu(m)}{m}\log\Delta(z_\star^m,u^m),
\]
where the series converges at an exponential rate.

\begin{corollary}[Real-algebraic reduction of \texorpdfstring{$\Delta(z_\star^m,u^m)$}{Delta(z*^m,u^m)} on the unit circle via completion]\label{cor:sync-kernel-weighted-logM-completed-real}
Take \(u=e^{i\theta}\) and let \(r=e^{i\theta/2}\), \(s=r+r^{-1}=2\cos(\theta/2)\in[-2,2]\). Write \(w_\star(s)\) for the smallest-modulus root of \(\widehat\Delta(w,s)=0\); then \(z_\star(u)=w_\star(s)/r\).
For any \(m\ge 1\), let
\[
S_m(s):=r^m+r^{-m}=2\cos(m\theta/2)=:C_m(s)\in\ZZ[s]
\]
(\(C_m\) is the Chebyshev--Adams operation; see Corollary~\ref{cor:primitive-completion-hatp}).
Then the following identity holds:
\[
\boxed{\ \Delta(z_\star(u)^m,u^m)=\widehat\Delta\!\bigl(w_\star(s)^m,\ S_m(s)\bigr)\ }.
\]
Therefore, on the unit-circle twist, every term of the \(\Delta\)-only series for \(\log\mathfrak{M}(\theta)\) can be evaluated in the real domain: the only inputs are the real parameter \(s\in[-2,2]\), the real number \(w_\star(s)\), and its powers.
\end{corollary}

\subsubsection{The First 10 Weighted Primitive Spectral Polynomials}\label{app:weighted-primitive-table}
The first 10 terms $p_n(u)=\sum_{\tau\ \mathrm{primitive},\,|\tau|=n}u^{S_\tau\varphi}$ are computed rigorously from
\(
p_n(u)=\frac{1}{n}\sum_{d\mid n}\mu(d)\,a_{n/d}(u^d)
\)
using the trace polynomials \(a_n(u)=\Tr(B(u)^n)\) of the matrix \(B(u)\); the one-click generated output is as follows (script \texttt{scripts/exp\_sync\_kernel\_weighted\_primitive\_completed.py}):

\input{sections/generated/eq_sync_kernel_weighted_primitive_pn_first10_en}
For $n\ge 2$, the $u^0$ and $u^n$ terms vanish, and the coefficients satisfy $p_{n,k}=p_{n,n-k}$.

\begin{corollary}[Completion of the primitive polynomials: \texorpdfstring{$\widehat p_n(s)\in\ZZ[s]$}{hat p in Z[s]} and the Chebyshev--Adams operation]\label{cor:primitive-completion-hatp}
Take
\[
u=r^2,\qquad s:=r+r^{-1},
\]
and for \(n\ge 1\) define the completed primitive polynomial
\[
\widehat p_n(s):=r^{-n}p_n(r^2).
\]
Then for all \(n\ge 1\) one has \(\widehat p_n(s)\in\ZZ[s]\). Moreover:
\begin{enumerate}
  \item If \(n\) is even, then \(\widehat p_n(s)\) is an even polynomial in \(s\);
  \item If \(n\) is odd, then \(\widehat p_n(s)\) is an odd polynomial in \(s\).
\end{enumerate}
Furthermore, for any integer \(d\ge 1\) define
\[
C_d(s):=r^d+r^{-d}\in\ZZ[s],
\qquad
C_d(s)=2T_d(s/2)
\]
(\(T_d\) is the Chebyshev polynomial of the first kind). Then the power substitution in the completed coordinate satisfies the following strict ``angle-doubling law'':
\[
r^{-dn}\,p_n(r^{2d})=\widehat p_n\!\bigl(C_d(s)\bigr),
\]
that is, the substitution \(u\mapsto u^d\) in the Witt/M\"obius inversion becomes, after completion, the Chebyshev--Adams operation \(s\mapsto C_d(s)\).
\end{corollary}

\begin{corollary}[Strict parity balance for odd lengths: \texorpdfstring{$p_n(-1)=0$}{p_n(-1)=0}]\label{cor:primitive-odd-parity-balance}
If \(n\ge 3\) is odd, then by the divisibility law \(u(u+1)\mid p_n(u)\) of Corollary~\ref{cor:primitive-odd-div},
\[
\boxed{\ p_n(-1)=\sum_{k\ge 0}(-1)^k p_{n,k}=0\ }.
\]
Hence, for odd-length primitive orbits, the counts on the two ``energy parity'' sides are strictly equal (in the signed-weight sense).
\end{corollary}
\begin{proof}
By the coefficient mirror symmetry \(p_{n,k}=p_{n,n-k}\) (see above),
\(
p_n(u)=u^n p_n(1/u)
\).
Setting \(u=r^2\) gives
\(
p_n(r^2)=r^{2n}p_n(r^{-2})
\),
so that \(r^{-n}p_n(r^2)\) is invariant under \(r\mapsto r^{-1}\). By the invariant subring identity
\(
\ZZ[r,r^{-1}]^{\,r\mapsto r^{-1}}=\ZZ[r+r^{-1}]=\ZZ[s]
\)
one obtains \(\widehat p_n(s)\in\ZZ[s]\).
\par
The parity assertion follows from the exponent parity: \(p_n(r^2)\) contains only powers \(r^{2k}\); after multiplication by \(r^{-n}\), the power indices share the same parity as \(n\). Since the even/odd decomposition of \(\ZZ[s]\) under \(s\mapsto -s\) agrees with the power parity under \(r\mapsto -r\), the stated parity of \(\widehat p_n\) follows.
\par
Finally, setting \(s=r+r^{-1}\) gives
\(
C_d(s)=r^d+r^{-d}
\),
which satisfies the recurrence \(C_{d+1}=sC_d-C_{d-1}\) and hence \(C_d\in\ZZ[s]\), equivalent to the Chebyshev polynomial relation \(C_d(s)=2T_d(s/2)\). Substituting \(u=r^2\) and \(u^d=r^{2d}\) yields
\[
r^{-dn}p_n(r^{2d})=\widehat p_n(r^d+r^{-d})=\widehat p_n(C_d(s)),
\]
which is the stated angle-doubling law. \qed
\end{proof}

\begin{corollary}[Completed Chebyshev--Witt transform: \texorpdfstring{$\widehat a\leftrightarrow \widehat p$}{hat a <-> hat p}]\label{cor:chebyshev-witt-completed}
In the same completed coordinate \(u=r^2,s=r+r^{-1}\), define the completed ghost (trace) polynomial
\[
\widehat a_n(s):=r^{-n}a_n(r^2)=r^{-n}\Tr(B(r^2)^n)\in\ZZ[s].
\]
Then for all \(n\ge 1\) the following strict identities hold:
\[
\boxed{\
\widehat p_n(s)=\frac1n\sum_{d\mid n}\mu(d)\,\widehat a_{n/d}\!\bigl(C_d(s)\bigr),
\qquad
\widehat a_n(s)=\sum_{d\mid n}d\,\widehat p_d\!\bigl(C_{n/d}(s)\bigr)
\ }.
\]
\end{corollary}
\begin{proof}
Apply the Witt/M\"obius inversion for \(u=r^2\),
\(
p_n(u)=\frac1n\sum_{d\mid n}\mu(d)a_{n/d}(u^d)
\),
multiply both sides by \(r^{-n}\), and use the completed angle-doubling law
\(
r^{-dn}f(r^{2d})=\widehat f(C_d(s))
\)
(applied to \(f=a_{n/d}\) and \(f=p_d\) respectively) to obtain the first identity; the second follows similarly from
\(
a_n(u)=\sum_{d\mid n} d\,p_d(u^{n/d})
\)
upon completion.
\end{proof}

\begin{corollary}[Prime-length congruence: \texorpdfstring{$\widehat a_p(s)\equiv C_p(s)\ (\mathrm{mod}\ p)$}{hat a_p congruent Cp mod p}]\label{cor:completed-prime-congruence}
Retaining the completed notation of Corollary~\ref{cor:chebyshev-witt-completed}, if \(p\) is prime then
\[
\boxed{\ \widehat a_p(s)\equiv C_p(s)\pmod p\ }.
\]
\end{corollary}
\begin{proof}
For a prime \(p\), from \(\widehat p_p(s)=\frac1p\bigl(\widehat a_p(s)-\widehat a_1(C_p(s))\bigr)\) (the first identity of Corollary~\ref{cor:chebyshev-witt-completed}, with only the divisors \(d=1,p\)) one sees that
\[
\widehat a_p(s)-\widehat a_1(C_p(s))\in p\,\ZZ[s].
\]
For this kernel, \(a_1(u)=\Tr(B(u))=p_1(u)=1+u\), so
\(
\widehat a_1(s)=r^{-1}(1+r^2)=r+r^{-1}=s
\), and hence \(\widehat a_1(C_p(s))=C_p(s)\).
Substituting back yields the congruence. \qed
\end{proof}

\begin{corollary}[Teichm\"uller-type congruence: \texorpdfstring{$\widehat a_p(s)\equiv s^p\ (\mathrm{mod}\ p)$}{hat a_p congruent s^p mod p}]\label{cor:completed-teichmueller-congruence}
Following the notation of Corollary~\ref{cor:completed-prime-congruence}, if \(p\) is a prime then the following congruence holds in \(\ZZ[s]\):
\[
\boxed{\ \widehat a_p(s)\equiv s^p\pmod p\ }.
\]
\end{corollary}
\begin{proof}
By Corollary~\ref{cor:completed-prime-congruence}, \(\widehat a_p(s)\equiv C_p(s)\pmod p\).
On the other hand, in \(\FF_p[r,r^{-1}]\) one has
\[
s^p=(r+r^{-1})^p\equiv r^p+r^{-p}=C_p(s)\pmod p,
\]
where the congruence follows from \(\binom{p}{k}\equiv 0\pmod p\) (\(1\le k\le p-1\)).
Thus \(C_p(s)\equiv s^p\pmod p\); combining these gives the conclusion. \qed
\end{proof}

\begin{corollary}[Chebyshev--Frobenius collapse: Frobenius rigidity of the completed Dwork congruence tower at \texorpdfstring{$\bmod p$}{mod p}]\label{cor:completed-dwork-frobenius-tower}
Let \(p\) be a prime and \(k\ge 1\). Under the completed-ghost notation, the following completed Dwork congruence holds:
\[
\boxed{\ \widehat a_{p^k}(s)\equiv \widehat a_{p^{k-1}}\!\bigl(C_p(s)\bigr)\pmod{p^k}\qquad\text{in }\ZZ[s]\ }
\]
(Theorem~\ref{thm:discussion-completed-dwork-chebyshev}).
Furthermore, at the \(\bmod p\) level this collapses to a pure Frobenius iteration:
\[
\boxed{\ \widehat a_{p^k}(s)\equiv \widehat a_{p^{k-1}}(s^p)\pmod p\qquad\text{in }\ZZ[s]\ },
\qquad
\boxed{\ \widehat a_{p^k}(s)\equiv s^{p^k}\pmod p\qquad\text{in }\ZZ[s]\ }.
\]
\end{corollary}
\begin{proof}
The first identity is a restatement of Theorem~\ref{thm:discussion-completed-dwork-chebyshev}.
For the second identity, reduce the first modulo \(\FF_p[s]\) and use the identity from the proof of Corollary~\ref{cor:completed-teichmueller-congruence},
\(
C_p(s)\equiv s^p\ (\mathrm{mod}\ p)
\),
to substitute the argument, yielding
\(
\widehat a_{p^k}(s)\equiv \widehat a_{p^{k-1}}(s^p)\ (\mathrm{mod}\ p)
\).
For the third identity, induct on \(k\): the base case \(k=1\) is Corollary~\ref{cor:completed-teichmueller-congruence};
if \(\widehat a_{p^{k-1}}(s)\equiv s^{p^{k-1}}\ (\mathrm{mod}\ p)\), then substituting \(s^p\) gives
\(
\widehat a_{p^{k-1}}(s^p)\equiv s^{p^k}\ (\mathrm{mod}\ p)
\), and combining with the second identity yields the conclusion. \qed
\end{proof}

\input{sections/generated/tab_sync_kernel_weighted_primitive_completed_en}
